\section{Introduction}\label{sec:intro}

In \emph{correlation clustering}~\cite{BansalBC04} the input is a complete
graph whose pairs are labelled \emph{similar} ($+$) or \emph{dissimilar}
($-$), and the task is to partition the vertices so that the number of
disagreements, $+$ pairs between clusters plus $-$ pairs inside clusters, is
minimum.  The number of clusters is free.  Equivalently, given the positive
graph $G^+=(V,\Ep)$, one adds or deletes the fewest edges that turn $G^+$ into
a disjoint union of cliques (\emph{cluster editing}).  The problem is
APX-hard~\cite{CharikarGW05}.  It arises in entity resolution, community
detection and the consolidation of inconsistent pairwise information.

Two lines of work have progressed separately.  Approximation ratios for
complete graphs have fallen from $3$ for Pivot~\cite{AilonCN08} through
$2.06$ by rounding the metric LP~\cite{ChawlaMSY15} to $1.485$ with the
cluster LP~\cite{CaoCLLNV24} and below~\cite{GarciaSorianoS26}.  In practice,
engineered heuristics such as KaPoCE~\cite{BlasiusEtAl21kapoceHeur}, the
winner of the PACE~2021 challenge~\cite{KellerhalsKNZ21pace}, and multilevel
memetic solvers~\cite{hausberger2025scc} find very good solutions on graphs
with millions of vertices.  What is missing on such graphs is a
\emph{certificate}: a lower bound on the optimum, checkable independently of
the solver, that tells how far a given clustering can be from optimal.  Exact
branch-and-bound solves instances with a few hundred
vertices~\cite{BlasiusEtAl22sea}, and its strong packing bounds are
implemented on a dense $n\times n$ matrix.  Solvers for the metric LP reach
about $10^4$ vertices~\cite{veldt2019metric,ruggles2020parallel,
SonthaliaG22project}.  The combinatorial bounds that do scale, packings of
bad triangles and the LP relaxations based on strong triadic
closure~\cite{veldt2022stc}, certify solutions only within a factor of about
two.

This paper is about the lower-bound side.  We compute lower bounds for
correlation clustering on sparse graphs with up to
$\numMaxCertN{}$ vertices and $\numMaxCertM{}$ edges, export them as
certificates that an independent program checks in exact arithmetic against
the raw instance file, and use them to certify the solutions of our own and of
other solvers.  The primal methods are secondary; they are included because a
bound is only as useful as the solutions it is compared with, and because the
certificate yields a decomposition of the optimality gap that guides a local
search.

\begin{figure}[!tbp]
\centering
\resizebox{\textwidth}{!}{% Overview of the two solvers and how their parts exchange information.
\begin{tikzpicture}[
  font=\footnotesize, >=Stealth,
  box/.style={draw=black!55, rounded corners=2pt, align=center, inner sep=2.5pt,
              minimum height=9mm, text width=30mm, fill=white},
  primal/.style={box, draw=pxblue, fill=pxblue!7},
  dual/.style={box, draw=pxorange, fill=pxorange!7},
  check/.style={box, draw=pxaqua, fill=pxaqua!9},
  outp/.style={box, text width=22mm},
  lab/.style={font=\scriptsize, text=black!70, align=center, fill=white, inner sep=1pt},
  arr/.style={->, draw=black!60, line width=0.6pt},
]
\node[box, text width=16mm] (in) at (-0.2,0) {input graph\\$G^+$};
\node[primal] (pivot) at (3.0,0) {Pivot, then\\iterated flipping};
\node[primal] (lns)   at (6.9,0) {gap-guided LNS\\\scriptsize exact MIPs on small blocks};
\node[primal] (px)    at (6.9,1.75) {PXMem\\\scriptsize contraction, crossover};
\node[outp]    (C)     at (10.9,0.88) {clustering $\mathcal C$\\\scriptsize$\mathbb E[\cost]\le3\,\OPT$};
\node[box, text width=34mm]    (gap)   at (6.9,-1.8) {gap decomposition\\\scriptsize$\cost(\mathcal C)-\LB=\sum g\ge0$};
\node[dual]   (sup)   at (3.2,-3.6) {distance-two\\support $P$};
\node[dual]   (bd)    at (6.9,-3.6) {block dual ascent\\\scriptsize anytime bound $\LB(y)$};
\node[check]  (chk)   at (10.9,-3.6) {independent checker\\\scriptsize parses the raw file};
\node[outp]    (lb)    at (14.0,-3.6) {certified\\$\lceil\LB\rceil\le\OPT$};
\node[outp]    (ratio) at (14.0,-1.8) {ratio bound\\$\cost(\mathcal C)/\lceil\LB\rceil$};

\draw[arr] (in) -- (pivot);
\draw[arr] (in.south) |- (sup.west);
\draw[arr] (pivot) -- (lns);
\draw[arr] (pivot.north) |- (px.west);
\draw[arr] (lns.east) -- ++(0.35,0) |- ([yshift=-1.5mm]C.west);
\draw[arr] (px.east) -- ++(0.35,0) |- ([yshift=1.5mm]C.west);
\draw[arr] (sup) -- (bd);
\draw[arr] (bd) -- node[lab, right] {rows, $y$} (gap);
\draw[arr] ([xshift=-5mm]lns.south) -- node[lab, left] {$\mathcal C$} ([xshift=-5mm]gap.north);
\draw[arr] ([xshift=5mm]gap.north) -- node[lab, right] {where to search} ([xshift=5mm]lns.south);
\draw[arr] (bd) -- node[lab, above] {certificate} (chk);
\draw[arr] (chk) -- (lb);
\draw[arr] (lb) -- (ratio);
\draw[arr] (C.east) -| (ratio.north);

\node[font=\scriptsize\bfseries, text=pxblue, anchor=east] at (3.0,1.75) {primal side};
\node[font=\scriptsize\bfseries, text=pxorange, anchor=west] at (1.9,-4.65) {dual side};
\end{tikzpicture}
}
\caption{Overview.  The dual side (orange) builds the distance-two support
and a lower bound by block dual ascent from a sparse star packing; the
certificate is re-checked by an independent program in exact arithmetic
against the raw instance.  The primal side (blue) computes clusterings: Pivot
and iterated flipping give a start, and either the gap-guided LNS of
CertiFlip or the memetic search PXMem improves it.  The gap decomposition
splits $\cost(\mathcal C)-\LB$ into non-negative local terms and tells the LNS
where to search.}
\label{fig:overview}
\end{figure}

\paragraph{Contributions.}
\begin{enumerate}
\item \textbf{Distance-two relaxation} (Section~\ref{sec:relax}).  Optimal
clusters have diameter at most two in $G^+$, so every pair at distance at
least three is separated in every optimal clustering.  We show that the
metric LP restricted to the pairs at distance at most two, with the triangle
inequalities through distance-three pairs kept as covering rows, is a
relaxation, and that each of its feasible points extends to a feasible point
of the full metric LP with the same cost (Theorem~\ref{thm:sparse}).  The
relaxation has $O(\sum_v\deg(v)^2)$ variables instead of $\binom n2$, and LP
rounding keeps its guarantees.  We add star and local subgraph inequalities,
which are valid for every clustering that separates the distance-three pairs.
This observation is simple; its value is that it makes the metric LP and its
strengthenings usable on graphs with $10^5$ to $10^6$ vertices.

\item \textbf{Anytime block-dual bound} (Section~\ref{sec:dual}).  We compute
a lower bound by dual block-coordinate ascent: each step solves the LP of one
vertex block exactly by cutting planes, with costs equal to the current
reduced costs.  The bound is valid after every step and never decreases, and
its state is one reduced cost per pair plus the rows with positive
multipliers (Theorem~\ref{thm:anytime}).  Any packing of induced stars or bad
triangles is a feasible start.  We compute one on the sparse support by a
ruin-and-recreate local search.  This sparse star packing is the scalable
counterpart of the root bound of the KaPoCE branch-and-bound.  On large
sparse graphs a long packing alone is stronger than a short packing followed
by block LPs in the same budget; the block LPs pay off on small and dense
instances (Sections~\ref{sec:exp-bounds} and~\ref{sec:exp-cert}).

\item \textbf{Exactly checked certificates} (Section~\ref{sec:cert}).  A
certificate lists the rows with positive multipliers.  A checker of about
\numCheckerLines{} lines that shares no code with the solver parses the raw
instance file, verifies that it is the instance the certificate was computed
for (SHA-256 of the file and of the canonical edge list), checks the validity
of every row by enumeration or by a closed form, and evaluates the bound in
integer arithmetic.  The statement that the checker evaluates is proved in
Lean~4 (Theorem~\ref{thm:anytime}(a)).  We archive every certificate behind
every reported bound.

\item \textbf{Gap decomposition and CertiFlip} (Section~\ref{sec:gap}).  For
any clustering, the difference between its cost and the dual bound splits
exactly into non-negative terms on pairs and rows (Theorem~\ref{thm:gap}).
Re-optimising a vertex set can gain at most the gap mass that touches it,
which certifies local optimality and selects neighbourhoods for a
large-neighbourhood search with exact sub-MIPs.  CertiFlip combines this with
a flipping local search after \citet{CohenAddadLPTYZ24} and returns a checked
bound with every solution.

\item \textbf{Computational study} (Section~\ref{sec:exp}).  On the 173
solved instances of the PACE~2021 exact track, the checked bound proves
optimality on \numPaceOptOurs{} and the root star packing of the KaPoCE
branch-and-bound on 79; averaged over instances the latter is stronger
(\numPaceMeanStar{} against \numPaceMeanOurs{} of the optimum), and the
comparison depends on density.  Our bound improves the published lower
bounds of two open instances.  On \numSnapCert{} SNAP graphs, \numSnapCertHeld{} of them from a
held-out set fixed in advance, the better of the block-dual bound and the
star packing alone certifies the best known solutions within
\numSnapGapLo{} to \numSnapGapHi{}, where scalable baselines certify
factors of \numBaseFacLo{} to \numBaseFacHi{}.  We compare with the sparse
star packing alone and with the metric LP on the support with triangle or
star rows, so that the contribution of each part is visible.
\end{enumerate}

A second primal solver, PXMem (Section~\ref{sec:primal}), recombines two
clusterings exactly block by block of their overlap.  We use it to measure how
tight the certificates are against strong solutions, and compare it with
KaPoCE under a protocol fixed in advance (Section~\ref{sec:exp-h2h}).  The
approximation guarantees we state for CertiFlip and PXMem, expected cost at
most $3\,\OPT$, are inherited from the Pivot start and serve only as a
safeguard; the quality of a solution is measured by its certificate.

\paragraph{Relation to certified computation.}  Certifying
algorithms~\cite{McConnellMNS11} return a witness with every output.  Safe
bounds from floating-point LP solutions~\cite{NeumaierS04safe,SteffyW13valid},
checkable certificates for mixed-integer programming~\cite{CheungGS17vipr}
and proof logging in SAT and pseudo-Boolean solving~\cite{WetzlerHH14drat,
GochtN21veripb} pursue the same goal.  Our certificates are simpler than
general MIP certificates because every row is local: its validity is decided
by enumerating the partitions of at most nine vertices or by a closed form.

\paragraph{Development and evaluation data.}  The methods were developed on
fifteen SNAP graphs and the PACE~2021 exact track.  Before the final
experiments we fixed a held-out set of twelve further SNAP graphs and the
analysis plan (Section~\ref{sec:exp-setup}); the PACE heuristic-track
instances that coincide with development graphs are reported separately.
Appendix~\ref{app:devlog} records which data informed which design decision.
